\paragraph{Supplement to Finding A51: weights and principal polarization do not force a small exponent.}
Fix $q=p^e$. For every integer $f\ge1$ there exists an ordinary,
principally polarized abelian variety $A_f/\mathbf F_q$ of dimension
$f$ such that
\[
 \begin{aligned}
 H_{A_f}(T)&=1-T^f+q^fT^{2f},\\
 A_f(\mathbf F_q)&\simeq\mathbf Z/q^f\mathbf Z,\\
 \dim_{\mathbf F_p}A_f(\mathbf F_q)[p]&=1,
 \qquad v_p\bigl(H_{A_f}(1)\bigr)=ef.
 \end{aligned}
 \tag{WV.1}
\]
Thus Weil weights, integrality, principal polarization, the displayed
cyclic Frobenius-polynomial shape, and one rational $p$-primary
generator do not by themselves give the small exponent required in
A47. This assertion concerns actual abelian varieties. No cyclic
deck action, Jacobian-cover realization, or compatibility in a
single geometric tower is asserted.

\begin{proof}
Put $Q=q^f$. Waterhouse's elliptic isogeny-class theorem gives
an elliptic curve $E/\mathbf F_Q$ with trace one: the trace is
prime to $p$ and $1<2\sqrt Q$. This is its ordinary case
\cite[Theorem~4.1(1), pp.~536--538; Corollary~4.4, p.~540]{Waterhouse1969AV}.
Consequently
\[
       H_E(U)=1-U+QU^2,
       \qquad \#E(\mathbf F_Q)=Q.
 \tag{WV.2}
\]
The finite group $E(\mathbf F_Q)$ is therefore a nonzero
$p$-group. An ordinary elliptic curve has exactly $p$ geometric
points killed by $p$, so this rational group has exactly one
$p$-primary generator. The elementary classification of finite
abelian groups makes it cyclic of order $Q$.

Take the Weil restriction
\[
               A_f=\operatorname{Res}_{\mathbf F_Q/\mathbf F_q}E.
 \tag{WV.3}
\]
It is an abelian variety of dimension $f$, geometrically the
product of the $f$ Frobenius conjugates of $E$, and is ordinary.
The canonical degree-one polarization of $E$ gives a principal
polarization of $A_f$: after base change it is the product of
those canonical polarizations. Preservation of the principal
polarization, including its descent, is the Weil-restriction
result in \cite[Proposition~1.22, pp.~27--28]{Diem2001Weil}.
The defining functor gives the equality of groups
\[
                  A_f(\mathbf F_q)=E(\mathbf F_Q).
 \tag{WV.4}
\]
On an auxiliary-adic Tate module, Frobenius has the induced
representation from $E$; the characteristic-polynomial formula
for Weil restriction is
\[
                  H_{A_f}(T)=H_E(T^f)
                              =1-T^f+q^fT^{2f}
 \tag{WV.5}
\]
\cite[Proposition~1.21, p.~26]{Diem2001Weil}.
Combining (WV.2), (WV.4), and (WV.5) proves every assertion
of (WV.1), in particular the exact exponent $q^f$ rather than
merely its order as an upper bound.

The Weil condition can also be read directly from the monic
polynomial
\[
                    X^{2f}-X^f+q^f.
 \tag{WV.6}
\]
If $\alpha$ is any root, then $\alpha^f$ is a root of
$Y^2-Y+q^f$. The latter two roots are complex conjugates of
absolute value $q^{f/2}$, so every $\alpha$ has absolute value
$\sqrt q$. The coefficients are integral and the reciprocal
identity for a dimension-$f$ abelian variety holds exactly.
Thus the large valuation is not produced by abandoning the
Weil constraints: the middle coefficient is simply $-1$,
and $1+(-1)$ cancels at $T=1$.
\end{proof}

For comparison with a putative transitive prime-power cover,
set $n=\ell^r$ and $f=\varphi(n)$, with $\ell$ odd and
$\ell\ne p$. These abelian varieties have the same polynomial
shape and rational generator number occurring in (TD.7), but
their exponents $M=q^{\varphi(n)}$ grow exponentially in $n$.
In particular they do not satisfy $n>(2M-1)^2$ for large $n$.
This defeats a deduction from only the listed numerical and
polarization constraints. It does not supply the integral
$\mathbf Z[\zeta_n]$ deck action, proper-conductor norm maps,
cover point-count congruences, or a fixed underlying genus-two
curve. The varieties $E$ may vary with $f$. Those additional
geometric and tower constraints remain possible sources of an
actual bound; no failure of such a bound for genuine curve
towers is proved here.
